\umaappendix{Especificaciones de OpenSpec}
\label{ap:especificaciones}

\section{Capacidades de OpenSpec}

La tabla~\ref{tab:capacidades} enumera las capacidades especificadas en \texttt{openspec/specs}. El ciclo de trabajo con estas especificaciones se describe en la sección~\ref{sec:sdd}.

{\small
\begin{longtable}{@{}p{0.38\textwidth} p{0.56\textwidth}@{}}
\caption{Capacidades especificadas con OpenSpec}\label{tab:capacidades}\\
\toprule
\textbf{Capacidad} & \textbf{Descripción} \\
\midrule
\endfirsthead
\toprule
\textbf{Capacidad} & \textbf{Descripción} \\
\midrule
\endhead
\texttt{admin-metrics-dashboard} & Página de métricas: consumo de IA de la organización, coste de ingesta y resultados de la Arena, según el rol. \\
\texttt{api-authorization} & Autorización en el servidor por rol, clave compartida en gRPC, errores sin detalles internos y arranque bloqueado sin secretos. \\
\texttt{blob-storage} & Almacenamiento compatible con S3 de los documentos (RustFS en desarrollo y Cloudflare R2 en producción). \\
\texttt{build-quality-gates} & Integración continua para todos los lenguajes, versiones centralizadas y despliegue solo tras la CI. \\
\texttt{bulk-reprocessing-pipeline} & Reprocesado en bloque con un \emph{worker} por pipeline. \\
\texttt{chat-history} & Historial de conversaciones por usuario y ayuntamiento. \\
\texttt{chat-ui-connection} & Conexión del chat al \emph{hub}: reconexión, varias conversaciones en paralelo e indicador de espera. \\
\texttt{chunking-embeddings} & Fragmentación, generación de \emph{embeddings} y almacenamiento de las pipelines planas. \\
\texttt{clean-architecture-boundaries} & Reglas de la arquitectura limpia y pruebas que las verifican. \\
\texttt{contexto-personalizado} & Temas y zonas de interés del usuario para adaptar las respuestas. \\
\texttt{evaluation-results} & Resultados puntuados de la evaluación automática guardados en la base de datos y mostrados en la página de métricas. \\
\texttt{go-scraper} & Servicio en Go que descarga, filtra, almacena y encola los documentos del BOE y del BOJA. \\
\texttt{hierarchical-retrieval} & Pipeline jerárquica: padres e hijos, enriquecimiento, expansión de consulta, búsqueda híbrida RRF y lenguaje claro. \\
\texttt{ingestion-administration} & Operaciones administrativas de ingesta detrás del servicio \texttt{IIngestionAdminService}. \\
\texttt{ingestion-consistency} & Consistencia entre PostgreSQL, Qdrant y el estado de los trabajos. \\
\texttt{ingestion-telemetry} & Métricas de coste de ingesta y tasa de victorias de la Arena. \\
\texttt{llm-provider-security} & Llamadas seguras a los proveedores de IA: TLS, clientes HTTP compartidos y claves solo para su proveedor. \\
\texttt{messaging-reliability} & Mensajería fiable con RabbitMQ: persistencia, confirmaciones, reintentos y cola de fallidos. \\
\texttt{mobile-navigation} & Menú lateral accesible en el móvil. \\
\texttt{multi-tenant-user-profiles} & Aislamiento por ayuntamiento con organizaciones de Auth0 y perfil de usuario. \\
\texttt{observability-grafana} & Telemetría con OpenTelemetry y Grafana, y métricas de IA sin contenido. \\
\texttt{procesamiento-historico} & Sugerencia de intereses a partir de las conversaciones recientes. \\
\texttt{qdrant-vector-store} & Qdrant como base de datos vectorial detrás del puerto \texttt{IVectorStore}. \\
\texttt{query-api-semantic-kernel} & Respuesta a preguntas con RAG sobre Semantic Kernel, con modelo por ayuntamiento y respaldo. \\
\texttt{question-arena} & Comparación a ciegas de las dos pipelines y registro de votos. \\
\texttt{rag-evaluation} & Evaluación automática y reproducible de las variantes RAG: corpus y preguntas fijos, RAGAS con un juez local, métricas de recuperación sin juez, latencia y coste de indexación. \\
\texttt{realtime-notifications} & Notificaciones de ingesta y del \emph{scraper} por SignalR, solo a quien puede verlas. \\
\texttt{scraper-dynamic-filtering} & Reglas de filtrado del \emph{scraper} en la base de datos y descargas bajo demanda por gRPC. \\
\texttt{session-timeout} & Cierre de sesión por inactividad. \\
\texttt{signalr-integration} & \emph{Hubs} de SignalR para el chat y las notificaciones. \\
\texttt{user-api-keys-vault} & Modelo de IA por ayuntamiento y claves de API cifradas. \\
\texttt{user-auth-auth0} & Inicio de sesión con Auth0 y validación de tokens en la API. \\
\bottomrule
\end{longtable}
}

La tabla recoge las capacidades de \texttt{openspec/specs} a 19 de noviembre de 2026.
% NOTA (Antonio): si se añaden capacidades nuevas antes de la entrega, revisar la tabla.
